\subsection{Performance Comparison of Deep Learning Approaches for EVSE Anomaly Detection}

In the following, it is important to summarize which machine learning and deep learning algorithms have shown promising performance in the field of anomaly detection for EVSE. Since the data involves time series and hardware-related measurements, the paper \cite{deeplearing_compare_EVSE}, titled \textit{Deep Learning-based Intrusion Detection System for Electric Vehicle Charging Station}, provides a valuable overview of suitable deep learning models in this context.\\
The authors compare a Deep Neural Network (DNN) with a Long Short-Term Memory (LSTM) model and conclude that the LSTM architecture achieves high accuracy—over 99\%—after only a few training epochs. This supports the later decision in this thesis to select the LSTM model for better performance and anomaly detection in EVSE environments.\\
\\
But the LSTM is only as strong as the significance of its features. Therefore, a good feature selection for the LSTM algorithm is essential to achieve an effective deep learning approach for anomaly detection. The paper by Kottilingal (2024) \cite{feature_selection} offers a systematic comparison of different feature selection methods in combination with an LSTM network for intrusion detection. The authors evaluate both the detection accuracy and the efficiency of the LSTM model when applying various feature selection strategies. The central finding is that reduced feature sets—such as those selected using tree-based methods—can achieve almost the same detection accuracy as larger sets, while significantly lowering model complexity and computational requirements. This aspect is particularly relevant for resource-constrained environments such as embedded EVSE systems.\\
In the context of EVSE, this means that implementing a practical and high-performance anomaly detection system requires not only choosing a suitable deep learning model (e.g., LSTM) but also effectively reducing the input dimensionality through feature selection. As already shown in \cite{EV_Security}, even with EVSE-specific data, selecting a smaller but relevant subset of hardware and kernel events can maintain detection accuracy while significantly reducing system load. The work by Kottilingal supports this approach with experimental results and provides practical recommendations for combining deep learning architectures with feature selection, making it a valuable reference for objective method comparisons in smart charging anomaly detection.
